\section{2026 年 11 月核验增补}

本节按论文首次提交时间排列。每项方法附有可核验的外部入选依据。

% METHOD: 2604.13333
\paragraph{SSD-GS：用分解的散射和阴影信息支持高斯重新照明}
2026.04~\cite{2604_13333}。SSD-GS 把原始 3DGS 的球谐外观换成漫反射、镜面反射、阴影与次表面散射四项，分别用解析或小型学习模块建模，再按阶段加入训练。可见性和散射在单一高斯的着色中承担不同作用，使改变光源后仍能计算阴影与柔和的散射项。

原文在合成和实拍逐灯点亮数据上对比重照明，并通过分项消融观察各项作用。这些结果依赖对应的照明采集和点光源协议；模型没有显式追踪多次反弹，光栅化及近似间接光不足以替代完整路径追踪。官方仓库提供训练、渲染与评价脚本。

\textbf{入选依据} \href{https://github.com/irisfreesiri/ssd-gs}{公开实现}已找到与论文对应的代码文件。

% METHOD: 2609.11223
\paragraph{Tri-DehazeGS：分离雾介质，并补偿被低透射率压弱的训练信号}
2026.09~\cite{2609_11223}。Tri-DehazeGS 将场景辐亮度与跨视图共享的三维雾介质分开表示。它同时补偿介质透射率降低后被压弱的重建梯度，使低可见区域能向场景表示提供更充分的约束。

论文在 RealX3D、Mip-NeRF 360 和 Fog-NeRF 比较，并在 RealX3D 上相对所列最强散射感知基线 3D-UIR 报告 PSNR 高 3.35 dB。重建仍依赖估计深度正则，细结构或低纹理处的深度过度平滑会导致边界模糊。

\textbf{入选依据} 论文署名单位包括 Meituan（\href{https://arxiv.org/abs/2609.11223}{论文原文}）。公开项目页未提供可确认的完整实现，按原文 Meituan 作者单位纳入。

% METHOD: 2609.39527
\paragraph{从多视图中分离并重建镜头眩光}
2026.09~\cite{2609_39527}。镜头眩光是随相机和光源改变的成像效应，将它写进静态场景高斯会把伪影误当作世界中的物体。作者先微调潜空间扩散模型得到去眩光视图，再用多视图共同拟合场景与眩光。眩光用位于场景前方二维平面上可移动的高斯表示，渲染时与场景层一起合成。

论文在 Flare7K++ 和自建 VFX 基准上评价二维去眩光，并另做多视图眩光重建和消融。二维去除指标与成像分解的质量要分别判断。镜头效应能否由前置平面模型表示，也取决于具体光学结构与观测条件。

\textbf{入选依据} 论文署名单位包括 Meta Reality Labs（\href{https://arxiv.org/abs/2609.39527}{论文原文}）。

% METHOD: 2611.24723
\paragraph{用分离的辐射高斯场渲染多波段星系}
2026.11~\cite{2611_24723}。该方法为恒星辐射、尘埃消光与尘埃辐射分别拟合高斯场，以物理基函数约束各场的光谱，并通过分阶段优化重建多波段视图。高斯表示承载各物理分量，不把同一像素的所有能量当成一种颜色系数。

评价使用 SKIRT-Auriga 和 TNG50 的模拟数据。其中 Auriga 子集为十个星系，每个以十七视图中的十四个训练、三个测试。作者指出，尘埃吸收的星光和其红外辐射未通过模型显式耦合，因此物理一致性仍有边界。

\textbf{入选依据} 论文署名单位包括 Baidu（\href{https://arxiv.org/abs/2611.24723}{论文原文}）。
